\appendixsection{Исходные тексты проекта}

Все исходные тексты, разработанные и использованные в рамках работы, размещены в
публичной GitHub-организации
\href{https://github.com/orgs/mai-diploma-kochkozharov/repositories}{\texttt{mai-diploma-kochkozharov}}
(\url{https://github.com/orgs/mai-diploma-kochkozharov/repositories}).
Организация содержит следующие репозитории.

\begin{itemize}
    \item \textbf{\href{https://github.com/mai-diploma-kochkozharov/paimon}{\texttt{paimon}}}
        (ветка \texttt{vortex-on-1.4.0-rc1}). Ответвление репозитория
        \texttt{apache/paimon} на базе версии~1.4. Содержит разработанные в
        рамках работы подмодули: \texttt{paimon-vortex-jni} (обёртки нативных
        методов и загрузчик библиотеки), \texttt{paimon-vortex-format}
        (реализация интерфейса \texttt{FileFormat} --- фабрики записи и чтения,
        \texttt{VortexFormatWriter}, \texttt{VortexRecordsReader}, конвертер
        предикатов \texttt{VortexPredicateConverter}), а также интегрированную
        сборку нативной библиотеки Vortex.

    \item \textbf{\href{https://github.com/mai-diploma-kochkozharov/paimon-flink-job}{\texttt{paimon-flink-job}}}
        (ветка \texttt{main}). Нагрузочные задания Apache Flink~+ Apache Paimon:
        задание потокового near-real-time-сценария
        \texttt{StreamingWriteBenchmarkJob}, задание пакетной записи
        \texttt{WriteBenchmarkJob}, исполнитель стандартизованного набора TPC-DS
        \texttt{TpcdsBenchmarkJob}, сценарий-обёртка многоформатного прогона и
        Maven-конфигурация сборки.

    \item \textbf{\href{https://github.com/mai-diploma-kochkozharov/flink-lang-deploy}{\texttt{flink-lang-deploy}}}
        (ветка \texttt{master}). Развёртывание кластера Apache Flink средствами
        Ansible: роли установки и конфигурации, плейбуки
        \texttt{start.yml}, \texttt{stop.yml}, \texttt{config-update.yml},
        \texttt{submit-job.yml}, \texttt{krb-renew.yml}, шаблоны конфигурации
        Apache Flink, инвентарь кластера.

    \item \textbf{\href{https://github.com/mai-diploma-kochkozharov/paimon-vortex-deploy}{\texttt{paimon-vortex-deploy}}}
        (ветка \texttt{glibc-2.28}). Сборка артефактов в контейнерах с
        кросс-компиляцией нативной библиотеки Vortex под целевую версию
        системной библиотеки C кластера; конфигурация Docker Compose локального
        окружения для функциональных smoke-тестов; Makefile-цели сборки и
        запуска заданий.

    \item \textbf{\href{https://github.com/mai-diploma-kochkozharov/vortex-jni-test}{\texttt{vortex-jni-test}}}
        (ветка \texttt{master}). Изолированный проект для тестирования слоя
        доступа к нативной библиотеке Vortex через Java Native Interface:
        проверка корректности экспорта данных в Apache Arrow без копирования,
        граничные случаи кодеков, отладка нативных вызовов вне основного
        конвейера Paimon.

    \item \textbf{\href{https://github.com/mai-diploma-kochkozharov/paimon-benchmark-results}{\texttt{paimon-benchmark-results}}}
        (ветка \texttt{master}). Результаты экспериментов и сценарии их обработки:
        файлы метрик TPC-DS и потокового сценария в формате CSV, сценарий разбора
        и агрегации \texttt{parse\_writebench.py} (поддерживает вычисление как
        среднего, так и медианы).

    \item \textbf{\href{https://github.com/mai-diploma-kochkozharov/diploma-latex-template}{\texttt{diploma-latex-template}}}
        (ветка \texttt{vkr-kochkozharov}). Исходные тексты настоящей выпускной
        квалификационной работы: \LaTeX-шаблон оформления, исходные
        \texttt{.tex}-файлы по разделам, конфигурация сборки (\texttt{latexmk},
        \texttt{xelatex}), исходники диаграмм в формате draw.io и их PDF-сборки.
\end{itemize}

Доступ ко всем репозиториям открытый; клонирование и сборка артефактов возможны
без авторизации.
